\documentclass[sigconf]{acmart}

\settopmatter{printacmref=false}
\renewcommand\footnotetextcopyrightpermission[1]{}
\pagestyle{plain}
\acmConference[NLP 2026]{NLP Course Project}{Spring 2026}{Northeastern University}
\acmYear{2026}

\usepackage{graphicx}
\usepackage{booktabs}
\usepackage{multirow}

\title{Ship of Theseus: Computational Forensics and the Paradox of Authorial Identity}

\author{Rohan Joshi}
\email{joshi.rohan@northeastern.edu}
\affiliation{%
  \institution{Northeastern University}
  \city{Seattle}
  \state{Washington}
  \country{USA}
}

\author{Aarav Surkatha}
\email{surkatha.aarav@northeastern.edu}
\affiliation{%
  \institution{Northeastern University}
  \city{Seattle}
  \state{Washington}
  \country{USA}
}

\begin{document}

\begin{abstract}
When you paraphrase a text multiple times using language models, the writing gradually stops sounding like the original author. This project investigates how authorial identity decays through iterative paraphrasing using the Ship of Theseus Corpus, which contains texts paraphrased across four generations ($T_0 \rightarrow T_3$) by seven paraphraser variants across four model families. We explore which linguistic features disappear fastest, identify the critical ``point of no return'' for authorship attribution, and detect model-specific fingerprints in paraphrased text. Our initial results confirm that semantic meaning is preserved (BERTScore 0.80--0.98) while lexical surface features are systematically replaced (BLEU as low as 0.002), with the sharpest decay occurring at the first iteration. Our work has implications for plagiarism detection, AI-generated content identification, and understanding the boundary between human and machine authorship.
\end{abstract}

\maketitle

\section{Introduction}

\subsection{Motivation}

The proliferation of large language models has made text paraphrasing trivially accessible. A student submitting coursework, a content farm scaling production, or a disinformation operator laundering provenance can each iteratively paraphrase a document through multiple LLMs, producing text that retains the original meaning while evading plagiarism detectors and authorship attribution systems. This creates urgent challenges for academic integrity, forensic linguistics, and AI transparency: if an LLM can rewrite text until no trace of the original author remains, traditional authorship verification collapses.

The threat is not hypothetical. Deployed detectors such as Turnitin report that a substantial fraction of student writing now contains AI-generated material at scale \cite{turnitin2024aiwriting}, and independent studies show that neural paraphrasers can evade state-of-the-art AI text detectors \cite{krishna2023dipper}, with iterative rewriting compounding the evasion effect. Emerging regulation such as the EU AI Act \cite{euaiact2024} expects AI-generated content to be labelled or otherwise flagged; those obligations are undermined whenever provenance signals can be stripped by rewriting. The consequences touch academic integrity offices deciding whether to escalate a suspected submission, newsrooms attempting to verify the origin of a leaked quote, and platforms enforcing synthetic-media policies. In each of these settings, a defender cannot simply ask \emph{was this text machine-generated?}: they must ask \emph{how much of the original author is still present in this version?} Yet most detection research still treats paraphrasing as a single-step perturbation, measuring whether a classifier can survive one rewrite. The more fundamental question---how authorial identity \emph{decays} across multiple generations of paraphrasing, which specific linguistic features erode first, and how quickly a usable forensic signal disappears---remains largely unexplored.

\subsection{Conceptual Framing}

We frame this problem through the Ship of Theseus paradox: if you replace every plank of a ship one by one, is the result still the same ship? Analogously, if an LLM replaces every word, restructures every sentence, and sheds every proper noun while preserving the original meaning, is the result still the original author's work? We map the philosophical framework onto measurable NLP signals: \emph{lexical planks} (individual words, measured by BLEU), \emph{structural hull} (grammatical patterns, measured by POS tag distributions), \emph{entity cargo} (factual anchors, measured by NER), and \emph{semantic meaning} (measured by BERTScore). This mapping allows us to quantify which ``planks'' are replaced first and at what rate.

\subsection{Research Gap}

Prior work on authorship attribution under adversarial conditions \cite{koppel2011authorship} demonstrates that paraphrasing degrades attribution accuracy, but does not quantify the \emph{rate} or \emph{order} of feature decay across iterations. Work on AI text detection \cite{mitchell2023detectgpt} focuses on binary classification (human vs.\ machine) rather than measuring how iterative rewriting erodes specific linguistic features at different rates. Neural paraphrase typologies \cite{wahle2023paraphrase} characterize the \emph{types} of changes paraphrasers make, but do not track how those changes accumulate across multiple generations.

The distinction between single-step and iterative paraphrasing is fundamental. Single-step studies ask a threshold question: ``can a detector survive one attack?'' Iterative studies ask a phase-transition question: ``at what depth does the authorial signal become unrecoverable?'' Traditional authorship attribution methods---from Mosteller and Wallace's Bayesian analysis of \emph{The Federalist Papers} to Burrows' Delta and its successors---were designed for static corpora where text is fixed. They provide no framework for measuring how stylometric features degrade under systematic, iterative transformation. Our work bridges this gap by treating authorial identity not as a binary property (present or absent) but as a continuously decaying signal whose components erode at measurably different rates.

Our work fills this gap by providing a multi-signal, multi-iteration forensic analysis that measures \emph{which} features decay, \emph{how fast}, and \emph{whether different paraphrasers leave identifiable signatures}. Treating identity as a decaying signal rather than a binary label lets us describe the forensic landscape at each iteration depth, and lets practitioners reason explicitly about which detectors remain trustworthy after one, two, or three rewrites.

\subsection{Research Questions}

We investigate three questions using the Ship of Theseus Paraphrased Corpus \cite{tripto2024ship}, each targeting a different aspect of the forensic problem: RQ1 characterizes \emph{what} decays and how fast, RQ2 identifies the \emph{threshold} at which identity is lost, and RQ3 asks whether the transformation tool itself can be identified---akin to determining which shipyard replaced the planks.

\begin{itemize}
    \item \textbf{RQ1:} Which linguistic features decay first---style (POS distributions, punctuation) or content (named entities, factual anchors)?
    \item \textbf{RQ2:} At what paraphrase iteration does authorship attribution effectively fail?
    \item \textbf{RQ3:} Do different paraphrasing models leave identifiable forensic fingerprints?
\end{itemize}

\textbf{Contributions.} We present a multi-signal forensic analysis combining lexical metrics (BLEU, ROUGE-L), semantic metrics (BERTScore), stylometric features (POS tag distributions), and content features (NER entity preservation) across seven paraphraser variants, seven diverse datasets, and three paraphrase iterations. We find that (1) different feature types exhibit \emph{layered erosion}, with lexical surface vanishing first, entities next, and grammatical structure last; (2) content features (named entities) are 5.8$\times$ more sensitive to paraphrasing intensity than stylistic features (POS distributions); (3) entities are predominantly \emph{dropped} rather than replaced by novel entities; and (4) each paraphraser occupies a distinct position in the style-vs-content preservation space, with paraphraser configuration mattering more than model architecture.

\section{Background}

\citet{stamatatos2009survey} provides a comprehensive survey of modern stylometric approaches that use lexical, syntactic, and structural features to identify text authorship. The survey examines both traditional statistical methods and machine learning approaches for capturing distinctive writing patterns. Although these methods achieve strong performance on controlled datasets, the robustness of stylometric features under systematic text transformations remains an open question.

\citet{koppel2011authorship} investigate authorship attribution in real-world scenarios where authors may deliberately obfuscate their writing style. They demonstrate that text transformations such as paraphrasing, translation round-trips, and stylistic rewriting significantly degrade the performance of standard authorship attribution systems. While their work establishes that obfuscation is possible, it does not systematically quantify the rate of stylometric decay across iterative transformations.

\citet{papineni2002bleu} introduce the BLEU metric, which provides an automatic method for evaluating machine translation quality through n-gram precision matching between generated and reference texts. This metric efficiently captures lexical overlap and has become the standard baseline for paraphrase evaluation in NLP. However, BLEU measures surface-level similarity rather than semantic equivalence, making it insufficient for capturing deeper authorial identity preservation.

\citet{lin2004rouge} presents ROUGE, a recall-oriented metric package originally designed for automatic summarization evaluation that measures n-gram overlap between system outputs and reference texts. ROUGE complements precision-focused metrics like BLEU by emphasizing content preservation. However, like BLEU, it remains a lexical metric and cannot distinguish between texts that share meaning but use entirely different words.

\citet{zhang2020bertscore} propose BERTScore, which evaluates text similarity using contextual BERT embeddings rather than surface n-gram overlap. BERTScore captures semantic equivalence more effectively than BLEU or ROUGE by comparing texts in embedding space. While this enables measurement of meaning preservation through paraphrasing, it was not designed to measure stylometric preservation, leaving a gap our work addresses by combining both lexical and semantic metrics to track the divergence between style decay and content preservation.
\section{Dataset}

We use the Ship of Theseus Paraphrased Corpus \cite{tripto2024ship}, a controlled benchmark for studying authorial identity under iterative paraphrasing.\footnote{\url{https://github.com/tripto03/Ship_of_theseus_paraphrased_copus}}

\subsection{Corpus Structure}

The corpus contains seven topically diverse datasets: \textbf{CMV} (Reddit ChangeMyView arguments), \textbf{ELI5} (Reddit explanations), \textbf{SCI\_GEN} (scientific text), \textbf{TLDR} (Reddit summaries), \textbf{WP} (WritingPrompts fiction), \textbf{XSum} (BBC news summaries), and \textbf{Yelp} (business reviews). These datasets span a wide range of registers (informal Yelp reviews to formal scientific text), discourse types (argumentative CMV to narrative WP), document lengths (47 words average for TLDR to 512 for WP), and entity densities (5.9 entities per document for Yelp to 18.7 for XSum). This diversity enables testing whether decay patterns generalize across domains. Each dataset includes texts from seven source authors: one human and six LLMs (OpenAI, LLAMA, Tsinghua, BigScience, PaLM, Eleuther-AI).

Each text is paraphrased sequentially three times to produce a four-stage chain: the original text $T_0$, first paraphrase $T_1$, second paraphrase $T_2$, and third paraphrase $T_3$, where each $T_n$ is generated by paraphrasing $T_{n-1}$.

\subsection{Paraphrasing Models}

Seven paraphraser variants are used across four model families:

\begin{itemize}
    \item \textbf{ChatGPT}: Full-text paraphrasing via OpenAI's GPT model.
    \item \textbf{PaLM2}: Full-text paraphrasing via Google's PaLM2.
    \item \textbf{Dipper} \cite{krishna2023dipper}: A controlled paraphraser with three diversity settings: default, low, and high.
    \item \textbf{Pegasus} \cite{zhang2020pegasus}: A sentence-level paraphraser with two modes: full (all sentences paraphrased) and slight (25\% of sentences paraphrased).
\end{itemize}

\subsection{Preprocessing}

After loading all seven datasets, we applied light preprocessing: removing PaLM2 markdown formatting artifacts, normalizing unicode whitespace, and filtering texts with fewer than 10 words. We preserved casing, punctuation patterns, and sentence structure since these are the stylometric features under investigation in RQ1. After deduplication on (document key, version), the processed corpus contains \textbf{66,400 rows}: 3,020 originals at $T_0$ and approximately 21,000 rows each at $T_1$, $T_2$, and $T_3$. Table~\ref{tab:corpus_stats} summarizes the corpus.

\subsection{Experimental Splits}

The experimental design is split-aware across two axes to avoid leakage. Along the \emph{tier} axis, RQ2 classifiers are trained exclusively on $T_0$ features and evaluated on $T_1$--$T_3$, so the model never sees a paraphrased copy of a document it was trained on. Along the \emph{document} axis, RQ3 paraphraser-identification uses \texttt{GroupShuffleSplit} with an 80/20 train--test ratio keyed on document identity, so every paraphrase of a given source document is confined to a single split. This prevents the classifier from memorising a document's $T_0$ signature and recognising it in paraphrased form at test time. Classifier hyperparameters are selected via 5-fold stratified cross-validation on the $T_0$ training set only; the held-out 20\% document split and the $T_1$--$T_3$ tiers are touched only at final evaluation.

\begin{table}[h]
\centering
\small
\begin{tabular}{lrrr}
\toprule
\textbf{Dataset} & \textbf{Paraphrased} & \textbf{Train} & \textbf{Versions} \\
\midrule
CMV & 5,626 & 1,666 & 22 \\
ELI5 & 10,447 & 3,181 & 22 \\
SCI\_GEN & 10,360 & 2,780 & 22 \\
TLDR & 8,404 & 2,559 & 22 \\
WP & 10,289 & 2,754 & 22 \\
XSum & 10,472 & 3,243 & 22 \\
Yelp & 10,802 & 3,344 & 22 \\
\midrule
\textbf{Total} & \textbf{66,400} & \textbf{19,527} & -- \\
\bottomrule
\end{tabular}
\caption{Corpus summary after preprocessing. Each dataset has 22 version strings (1 original + 7 paraphrasers $\times$ 3 iterations).}
\label{tab:corpus_stats}
\end{table}
\section{Experimentation and Results}

\subsection{Similarity Baselines}

Table~\ref{tab:baselines} presents BLEU, ROUGE-L, and BERTScore between $T_0$ and each paraphrased tier. Figure~\ref{fig:decay} visualizes the aggregate decay.

\begin{table}[h]
\centering
\small
\begin{tabular}{llccc}
\toprule
\textbf{Paraphraser} & \textbf{Metric} & \textbf{T1} & \textbf{T2} & \textbf{T3} \\
\midrule
\multirow{3}{*}{ChatGPT} & BLEU & .195 & .172 & .143 \\
 & ROUGE-L & .503 & .464 & .426 \\
 & BERTScore & .923 & .915 & .907 \\
\midrule
\multirow{3}{*}{PaLM2} & BLEU & .470 & .383 & .338 \\
 & ROUGE-L & .698 & .625 & .584 \\
 & BERTScore & .953 & .943 & .938 \\
\midrule
\multirow{3}{*}{Dipper} & BLEU & .168 & .093 & .061 \\
 & ROUGE-L & .377 & .317 & .257 \\
 & BERTScore & .905 & .888 & .877 \\
\midrule
\multirow{3}{*}{Dipper (High)} & BLEU & .010 & .004 & .002 \\
 & ROUGE-L & .130 & .107 & .091 \\
 & BERTScore & .806 & .804 & .803 \\
\midrule
\multirow{3}{*}{Dipper (Low)} & BLEU & .492 & .369 & .304 \\
 & ROUGE-L & .716 & .628 & .565 \\
 & BERTScore & .953 & .940 & .932 \\
\midrule
\multirow{3}{*}{Pegasus (Full)} & BLEU & .415 & .297 & .252 \\
 & ROUGE-L & .673 & .588 & .540 \\
 & BERTScore & .931 & .914 & .906 \\
\midrule
\multirow{3}{*}{Pegasus (Slight)} & BLEU & .841 & .728 & .638 \\
 & ROUGE-L & .919 & .857 & .804 \\
 & BERTScore & .980 & .967 & .956 \\
\bottomrule
\end{tabular}
\caption{Similarity $T_0$ vs.\ $T_n$. BLEU/ROUGE-L lexical; BERTScore semantic. Corpus means.}
\label{tab:baselines}
\end{table}

\begin{figure}[h]
\centering
\includegraphics[width=\columnwidth]{figures/baselines/aggregate_style_vs_content.png}
\Description{Line chart showing BERTScore staying above 0.80 while BLEU drops below 0.40 across T0 to T3.}
\caption{Aggregate style vs.\ content decay (RQ1). BERTScore stays above 0.80 while BLEU drops below 0.40 by $T_1$: meaning survives, surface is replaced.}
\label{fig:decay}
\end{figure}

Across all paraphrasers, BERTScore remains 0.80--0.98 while BLEU drops as low as 0.002: meaning is preserved while surface is replaced. The sharpest decay occurs at $T_1$, with only modest additional loss through $T_3$. Within-family variation (Dipper BLEU range: 0.002--0.492) exceeds between-family variation, which is preliminary evidence for RQ3. Exploratory analysis also shows monotonic text compression (279$\rightarrow$184 words) and punctuation erosion across tiers.

\subsection{POS Tag Distributions: Stylistic Drift (RQ1)}

Table~\ref{tab:pos_cosine} reports POS cosine similarity ($T_0$ vs.\ $T_1$) for all seven paraphrasers. The ranking mirrors the lexical ordering from \S5.1.

\begin{table}[h]
\centering
\small
\begin{tabular}{lcc}
\toprule
\textbf{Paraphraser} & \textbf{POS Cosine} & \textbf{Std Dev} \\
\midrule
Pegasus (Slight) & .9922 & .0396 \\
Dipper (Low) & .9840 & .0292 \\
PaLM2 & .9755 & .0385 \\
ChatGPT & .9686 & .0381 \\
Dipper (Default) & .9637 & .0424 \\
Pegasus (Full) & .9626 & .0834 \\
Dipper (High) & .8782 & .0906 \\
\bottomrule
\end{tabular}
\caption{POS cosine ($T_0$ vs.\ $T_1$). Higher = less syntactic restructuring. Even Dipper (High) retains 0.88---grammar is remarkably resilient.}
\label{tab:pos_cosine}
\end{table}

\textbf{Key finding:} Grammatical structure is resilient. Even Dipper (High), which pushes BLEU to 0.002, retains POS cosine 0.878. The spread within the Dipper family alone (0.984--0.878) exceeds the spread across all other model families, reinforcing that \emph{paraphraser configuration matters more than model architecture}.

\subsection{NER Entity Preservation: Content Erasure (RQ1)}

Table~\ref{tab:ner_metrics} presents entity-level metrics for all seven paraphrasers. Articles with zero $T_0$ entities (1,590 of 18,058, or 8.8\%) are excluded from means.

\begin{table}[h]
\centering
\small
\begin{tabular}{lccc}
\toprule
\textbf{Paraphraser} & \textbf{Jaccard} & \textbf{Recall} & \textbf{Precision} \\
\midrule
Pegasus (Slight) & .868 & .887 & .967 \\
PaLM2 & .590 & .658 & .798 \\
Dipper (Low) & .557 & .648 & .734 \\
ChatGPT & .552 & .629 & .766 \\
Pegasus (Full) & .517 & .557 & .846 \\
Dipper (Default) & .330 & .424 & .537 \\
Dipper (High) & .021 & .038 & .059 \\
\bottomrule
\end{tabular}
\caption{NER entity metrics ($T_0$ vs.\ $T_1$). Low Recall = dropping; low Precision = novel-entity introduction. Dipper (High) destroys 96\% of named entities.}
\label{tab:ner_metrics}
\end{table}

\textbf{Key finding:} Entities are far more fragile than grammar. Dipper (High) retains only 3.8\% of original entities (Recall 0.038) while its POS cosine stays at 0.878. This asymmetry answers RQ1: content decays faster than style.

\subsection{Controlled Experiments: The 5.8$\times$ Sensitivity Gap (RQ1)}

Using Dipper Low vs.\ High as a controlled comparison, content features (NER Recall delta: 0.610) are \textbf{5.8$\times$ more sensitive} to paraphrasing intensity than stylistic features (POS cosine delta: 0.104); Figure~\ref{fig:dipper_intensity} visualises the asymmetry. The Pegasus coverage experiment confirms that paraphrasers primarily \emph{shed} entities (Recall $-$0.329) rather than introduce novel ones (Precision $-$0.122).

\begin{figure}[h]
\centering
\includegraphics[width=\columnwidth]{figures/fingerprints/q3_dipper_intensity.png}
\Description{Two bar charts comparing Dipper Low vs High: POS cosine drops 0.104, NER Recall drops 0.610.}
\caption{Dipper intensity: content is 5.8$\times$ more sensitive than style (RQ1).}
\label{fig:dipper_intensity}
\end{figure}

\subsection{Paraphraser Fingerprints (RQ3)}

Figure~\ref{fig:fingerprints} plots each paraphraser in style-vs-content space; each occupies a distinct region, confirming identifiable forensic fingerprints.

\begin{figure}[h]
\centering
\includegraphics[width=\columnwidth]{figures/fingerprints/style_vs_content_scatter.png}
\Description{Scatter plot showing seven paraphrasers positioned by POS cosine (x-axis) and NER Jaccard (y-axis), with quadrant labels.}
\caption{Paraphrasers in POS cosine (style) vs.\ NER Jaccard (content) space (RQ3). Pegasus (Slight) preserves both; Dipper (High) destroys both; models with similar POS can diverge sharply on NER.}
\label{fig:fingerprints}
\end{figure}

Notably, Pegasus (Full) and Dipper (Default) share POS cosine ($\approx$0.963) but diverge on NER Precision (0.846 vs.\ 0.537): fingerprints are \emph{multi-dimensional}, and combining style and content signals produces separable signatures.

\subsection{Multi-Tier Decay: T0 Through T3 (RQ1)}

Extending to $T_2$ and $T_3$, Table~\ref{tab:multitier} reports POS cosine and NER Recall at every tier; Figure~\ref{fig:composite_decay} visualizes the aggregate pattern.\footnote{T1 values in Tables~\ref{tab:pos_cosine} and \ref{tab:ner_metrics} differ slightly from this table because the multi-tier pivots are aligned per-paraphraser across all four tiers, producing different sample populations (508--3,018 per paraphraser) than the 18,058-row all-paraphrasers pivot.}

\begin{table}[h]
\centering
\small
\begin{tabular}{lcccccc}
\toprule
 & \multicolumn{3}{c}{\textbf{POS Cosine}} & \multicolumn{3}{c}{\textbf{NER Recall}} \\
\cmidrule(lr){2-4} \cmidrule(lr){5-7}
\textbf{Paraphraser} & \textbf{T1} & \textbf{T2} & \textbf{T3} & \textbf{T1} & \textbf{T2} & \textbf{T3} \\
\midrule
Pegasus (Slight) & .993 & .986 & .980 & .892 & .807 & .736 \\
Dipper (Low) & .984 & .979 & .975 & .651 & .567 & .521 \\
PaLM2 & .978 & .974 & .972 & .715 & .672 & .635 \\
ChatGPT & .969 & .966 & .961 & .640 & .606 & .474 \\
Dipper (Default) & .965 & .956 & .948 & .422 & .326 & .272 \\
Pegasus (Full) & .964 & .948 & .938 & .559 & .460 & .418 \\
Dipper (High) & .884 & .863 & .843 & .043 & .018 & .012 \\
\bottomrule
\end{tabular}
\caption{POS cosine and NER Recall ($T_0$ vs.\ $T_n$) across tiers. POS decays gradually; NER Recall drops steeply (Dipper High: 1.2\% at $T_3$).}
\label{tab:multitier}
\end{table}

\begin{figure}[h]
\centering
\includegraphics[width=\columnwidth]{figures/baselines/feature_decay_layered.png}
\Description{Line chart showing six features decaying at different rates from T0 to T3: BERTScore and POS cosine stay above 0.90, NER Recall drops to 0.43, BLEU and ROUGE-L collapse below 0.30.}
\caption{Composite feature decay $T_0{\rightarrow}T_3$ aggregated over all paraphrasers (RQ1). Layered erosion: BERTScore and POS cosine stay above 0.90; NER Recall falls to 0.43; BLEU/ROUGE-L collapse below 0.30.}
\label{fig:composite_decay}
\end{figure}

Layered erosion holds at all depths: the ordering BERTScore $>$ POS cosine $>$ NER Recall $>$ ROUGE-L $>$ BLEU is preserved at every tier. Even at $T_3$, Dipper (High) retains POS cosine 0.843 while NER Recall collapses to 1.2\%. Dependency depth is nearly flat across tiers.

\subsection{Authorship Attribution Decay (RQ2)}

Following the protocol in \S4.3, Figure~\ref{fig:attribution_decay} shows F1 decay curves across six paraphrasers and three classifiers. For Logistic Regression, macro F1 drops below the 0.5 random baseline at $T_1$ under Dipper (High): a single aggressive iteration is enough to erase stylometric identity. Pegasus (Slight) retains the most signal, with ensemble methods staying above 0.5 through $T_3$.

\begin{figure}[h]
\centering
\includegraphics[width=\columnwidth]{figures/fingerprints/attribution_decay_binary.png}
\Description{Three-panel chart showing F1 decay from T0 to T3 for six paraphrasers across Logistic Regression, SVM, and Random Forest classifiers.}
\caption{Binary (Human vs.\ LLM) F1 decay across tiers (RQ2). Dashed gray line is the 0.5 random baseline. Dipper (High) drops below baseline at $T_1$; Pegasus (Slight) retains signal through $T_3$.}
\label{fig:attribution_decay}
\end{figure}

\subsection{Paraphraser Identification (RQ3)}

Following the protocol in \S4.4, SVM (RBF) achieves the highest macro F1 of \textbf{0.755} at 6-class paraphraser identification, confirming that each paraphraser leaves a distinct fingerprint in the delta feature space. Figure~\ref{fig:confusion}: Dipper (High) is most identifiable (93\% correct) while Dipper (Low) is most confusable with Dipper Default.

\begin{figure}[h]
\centering
\includegraphics[width=\columnwidth]{figures/fingerprints/paraphraser_confusion_matrix.png}
\Description{Confusion matrix showing 6-class paraphraser identification accuracy. Dipper High is 93 percent correct; Dipper Low is 57 percent correct.}
\caption{6-class paraphraser confusion matrix (SVM, macro F1 = 0.755, RQ3). Dipper (High) most distinctive; Pegasus variants well-separated from Dipper.}
\label{fig:confusion}
\end{figure}

Random Forest feature importance (Figure~\ref{fig:importance}) ranks NER Jaccard, word count ratio, and NER Precision as the top three; POS deltas contribute moderately, with AUX and NOUN shifts most informative.

\begin{figure}[h]
\centering
\includegraphics[width=\columnwidth]{figures/fingerprints/paraphraser_feature_importance.png}
\Description{Horizontal bar chart showing feature importance for paraphraser identification. NER Jaccard, word count ratio, and NER Precision are the top three features.}
\caption{Feature importance for paraphraser identification (Random Forest, RQ3). Entity-level features and text compression dominate POS shifts, consistent with the RQ1 finding that content decays faster than style.}
\label{fig:importance}
\end{figure}

\subsection{Domain-Specific Analysis}

Decay varies by domain (Figure~\ref{fig:domain_heatmap}). Entity-rich XSum (18.7 entities/doc) retains NER better (Recall 0.461) than Yelp (5.9 entities, 0.362); CMV has the most resilient POS (0.976); ELI5 attribution collapses nearest chance ($F_1$=0.542 at $T_1$).

\begin{figure}[h]
\centering
\includegraphics[width=\columnwidth]{figures/domain/domain_paraphraser_heatmap.png}
\Description{Two heatmaps showing POS cosine and NER Recall by dataset and paraphraser.}
\caption{Domain $\times$ paraphraser preservation heatmap (RQ1). Entity-rich domains retain more content; POS resilience is highest in argumentative text.}
\label{fig:domain_heatmap}
\end{figure}

Figure~\ref{fig:domain_attribution} shows domain-wise attribution decay: CMV retains the strongest signal (F1 0.706 at $T_1$), while ELI5 and Yelp collapse toward random by $T_1$. Forensic tools likely need domain-specific calibration.

\begin{figure}[h]
\centering
\includegraphics[width=\columnwidth]{figures/domain/domain_attribution_decay.png}
\Description{Line chart showing attribution F1 decay across T0-T3 for seven datasets. CMV and XSum retain the most signal; ELI5 and Yelp collapse toward random.}
\caption{Attribution F1 decay by domain (Dipper, Logistic Regression, RQ2). CMV and XSum retain signal longest; ELI5 and Yelp collapse to random baseline by $T_1$.}
\label{fig:domain_attribution}
\end{figure}

\subsection{Identity Trajectory in Feature Space}

Applying t-SNE to the 23-dim feature vectors at each tier for Dipper (3,010 articles), Figure~\ref{fig:tsne} shows $T_0$ documents occupying a diverse region while $T_3$ clusters into a homogeneous region---visual evidence of convergence toward a machine-normalised style.

\begin{figure}[h]
\centering
\includegraphics[width=\columnwidth]{figures/fingerprints/identity_trajectory_tsne.png}
\Description{Two t-SNE scatter plots: left colored by paraphrase tier showing T3 clusters more tightly than T0; right showing Human vs LLM centroids converging from T0 to T3.}
\caption{Identity trajectory t-SNE projection (RQ2). Left: colored by tier, showing progressive clustering. Right: Human vs.\ LLM centroids at $T_0$ (circles) and $T_3$ (crosses) converge under paraphrasing.}
\label{fig:tsne}
\end{figure}

To validate the visual impression quantitatively (t-SNE can mislead), we compute silhouette, $k$-NN purity ($k{=}10$), and Human$\leftrightarrow$LLM centroid distance at each tier. All three decline monotonically from $T_0$ to $T_3$ (Figure~\ref{fig:tsne_quant}): centroid distance drops 0.426$\rightarrow$0.287 and $k$-NN purity 0.881$\rightarrow$0.839. Identity is eroding in feature space, not merely compressing visually.

\begin{figure}[h]
\centering
\includegraphics[width=\columnwidth]{figures/fingerprints/identity_quantitative_metrics.png}
\Description{Three bar charts showing silhouette score, k-NN purity, and centroid distance decreasing from T0 to T3.}
\caption{Quantitative validation of identity erosion (RQ2). Silhouette, $k$-NN purity, and Human$\leftrightarrow$LLM centroid distance all decline $T_0{\rightarrow}T_3$.}
\label{fig:tsne_quant}
\end{figure}
\section{Plan for Next}

In the coming weeks, we will pursue three directions:

\textbf{Stylometric feature extraction (Weeks 7--8):} We will extract at least five structural features per document: POS tag frequency distributions, sentence length variance, dependency tree depth, punctuation patterns, and Type-Token Ratio. These will be normalized relative to $T_0$ and ranked by decay rate to determine which linguistic ``planks'' are replaced first (RQ1).

\textbf{Authorship attribution (Weeks 9--10):} We will train classifiers (Logistic Regression, SVM) exclusively on $T_0$ data and evaluate on $T_1$--$T_3$ to identify the ``point of no return'' where attribution accuracy collapses (RQ2). Training and testing will never use the same paraphrase level.

\textbf{Paraphraser identification (Weeks 11--12):} Building on the fingerprint evidence from our baseline analysis, we will train multi-class classifiers to distinguish which paraphraser was used. Evaluation will include confusion matrices and feature importance analysis (RQ3).

\bibliographystyle{ACM-Reference-Format}
\bibliography{references}

\end{document}